\documentclass{handout}
\usepackage{handoutSetup}\usepackage{handoutShortcuts}  
\usepackage[noheads,nolists,nomarkers]{endfloat}

\begin{document}
\handoutHeader

\begin{verbatimwrite}{\jobname.title}
Decentralizing the Ramsey/Cass-Koopmans Model
\end{verbatimwrite}

\handoutNameMake

\newboolean{NumericalSolution}
\setboolean{NumericalSolution}{true}
%\setboolean{NumericalSolution}{false}


%\renewcommand{\cite}{\citeyearpar}

This handout shows that under certain very special conditions the
behavior of an economy composed of distinct individual households will
replicate the social planner's solution to the Ramsey/Cass-Koopmans
model.

\section{The Consumer's Problem}

Consider first the problem of an individual infinitely lived consumer
indexed by $i$ who has some predetermined set of expectations for 
how the aggregate net interest rate $\rfree_{t}$ and wage rate $\Wage_{t}$ will
evolve.

At date $t$ household $i$ owns some capital $\kap_{t,i}$, and can in principle also
borrow; designate the net debt of household $i$ in period $t$ as
$\debt_{t,i}$; since we will be examining a perfect foresight solution with
perfect capital markets, the interest rate on debt must match the 
rate of return on assets.  This means that all that really matters is
the household's total net asset position,
\opt{MarginNotes}{\marginpar{\tiny Each household can
    `issue bonds' which is like borrowing; you promise to pay interest
    to the bondholder.}}
\begin{equation}\begin{gathered}\begin{aligned}
        \wNet_{t,i} & =  \kap_{t,i} - \debt_{t,i}.
\end{aligned}\end{gathered}\end{equation}

Each household is endowed with one unit of labor, which it 
supplies exogenously, earning a wage rate $\Wage_{t,i}$.  \opt{MarginNotes}{\marginpar{\tiny We will determine below how $\Wage_{t}$ is determined; take it as exogenous here.}} 
Each household solves:
\begin{equation}
        \max \int_{0}^{\infty} \uFunc(\cons_{t,i}) e^{-\timeRate t} dt
\end{equation}
subject to the budget constraint
\begin{equation}
        \dot{\wNet}_{t,i} = \Wage_{t} + \rfree_{t} \wNet_{t,i} - \cons_{t,i}
.
\end{equation}

Integrating the household's dynamic budget constraint and assuming 
a no-Ponzi-game transversality condition yields the 
intertemporal budget constraint, which says that the present 
discounted value of consumption must match the PDV of labor
income plus the current stock of net wealth:
\begin{equation}\begin{gathered}\begin{aligned}
        \mathbb{P}_{0,i}(\cons) & =  \mathbb{P}_{0,i}(\Wage)+\wNet_{0,i}.
\end{aligned}\end{gathered}\end{equation}  

The formulas for these PDV's are a bit awkward because they must
take account of the fact that interest rates are varying over time.
To make the formulas a bit simpler, define the compound interest factor
\begin{equation}\begin{gathered}\begin{aligned}
        \RCpnd_{t}^{-1} & =  \exp(-\int_{0}^{t} \rfree_{\tau} d\tau),
\end{aligned}\end{gathered}\end{equation}
which is simply the compound interest term needed to convert a value 
at date $t$ to its PDV as of time 0.

With this definition in hand we can write the IBC as
\begin{equation}\begin{gathered}\begin{aligned}
        \int_{0}^{\infty} \cons_{t,i} \RCpnd^{-1}_{t} dt & =  h_{0,i}+ \wNet_{0,i}
\end{aligned}\end{gathered}\end{equation}
where $h_{0,i}$ is human wealth,
\begin{equation}\begin{gathered}\begin{aligned}
        h_{0,i} & =  \int_{0}^{\infty} \Wage_{t} \RCpnd^{-1}_{t} dt. \label{eq:humwealth}
\end{aligned}\end{gathered}\end{equation}

Each household solves the standard optimization problem taking 
the future paths of wages and interest rates as {\it given}.  Thus
the Hamiltonian\footnote{See \handoutG{RamseyCassKoopmans} for the discounted
Hamiltonian optimality conditions and \handoutG{HamiltonianVSDiscrete} for the intuition
of the logic behind the Hamiltonian.} is
\begin{equation}
        \Ham(\cons_{t,i},\wNet_{t,i},\lambda_{t,i}) = \uFunc(\cons_{t,i}) + (\rfree_{t}\wNet_{t,i}+\Wage_{t}-\cons_{t,i})\lambda_{t,i}
\end{equation}
which implies that the first optimality condition is the usual $\uP(\cons_{t,i}) = \lambda_{t,i}$.
The second optimality condition is
\begin{equation}\begin{gathered}\begin{aligned}
        \dot{\lambda}_{t,i} & =  \timeRate \lambda_{t,i} - (\partial \Ham/\partial \wNet) \\
        \dot{\lambda}_{t,i} & =  \timeRate \lambda_{t,i} - \lambda_{t,i} \rfree_{t} \\
        \dot{\lambda}_{t,i}/\lambda_{t,i} & =   (\timeRate - \rfree_{t})
\end{aligned}\end{gathered}\end{equation}
leading eventually to the usual first order condition for consumption:
\begin{equation}\begin{gathered}\begin{aligned}
        \dot{\cons}_{t,i}/\cons_{t,i} & =  \CRRA^{-1}(\rfree_{t}-\timeRate) \label{eq:cdotindiv}.
\end{aligned}\end{gathered}\end{equation}

Note (for future use) that the RHS of this equation does not contain
any components that are idiosyncratic: The consumption growth rate will
be identical for every household.  The same is true of the expression
for human wealth, equation \eqref{eq:humwealth}.

\section{The Firm's Problem}

Now we assume that there are many perfectly competitive small firms \opt{MarginNotes}{\marginpar{\tiny Because of the perfect competition and CRS assumptions, it doesn't matter exactly how many firms there are; firms are `atomistic' like households; like HH's, they take $\Wage_{t}$ and $\rfree_{t}$ as given.}} indexed by $j$
in this economy, each of which has a production function identical to
the aggregate Cobb-Douglas production function.  Perfect competition 
implies that individual firms take the interest
rate $\hat{\rfree}_{t}$ and wage rate $\hat{\Wage}_{t}$ to be exogenous.  Hence
firms solve \opt{MarginNotes}{\marginpar{\tiny Note that we are back to the Hall-Jorgensen model with no adjustment costs, so the firm problem has no important intertemporal dimension.}}
\begin{equation}
        \max_{\{K_{t,j},L_{t,j}\}} \FFunc(K_{t,j},L_{t,j}) - \hat{\Wage}_{t}L_{t,j} - \hat{\rfree}_{t}K_{t,j}
\end{equation}
where $\hat{\rfree}_{t}$ and $\hat{\Wage}_{t}$ are the rental rates for a unit 
of capital and a unit of labor for one period.  Note that, dividing by 
$L_{t,j}$, this is equivalent to
\begin{equation}
        \max_{\{ \kap_{t,j} \}} \underbrace{\kap_{t,j}^{\alpha}}_{\fFunc(\kap_{t,j})} - \hat{\Wage}_{t} - \hat{\rfree}_{t} \kap_{t,j}.
\end{equation}

The first order condition for this problem implies that
\begin{equation}\begin{gathered}\begin{aligned}
        \fFunc^{\prime}(\kap_{t,j}) & =  \hat{\rfree}_{t}.
\end{aligned}\end{gathered}\end{equation}

Under perfect competition firms must make zero profits in equilibrium,
which means, by fact \EulersTheorem, that:
\begin{equation}\begin{gathered}\begin{aligned}
        \fFunc(\kap_{t,j}) & =  \hat{\Wage}_{t}+\hat{\rfree}_{t}\kap_{t,j}.
\end{aligned}\end{gathered}\end{equation}

\section{Equilibrium At a Point in Time}

Thus far, we have solved the consumer's and the firm's problems 
from the standpoint of atomistic individuals.  It is now time to 
consider the behavior of an aggregate economy composed of consumers
and firms like these.

We assume that the population of households and firms is distributed
along the unit interval and the population masses sum to one, 
as per \handoutF{Aggregation}.  Thus, aggregate assets at time $t$
can be defined as the sum of the assets of all the individuals in the
economy at time $t$,
\begin{equation}\begin{gathered}\begin{aligned}
  \WNet_{t} & =  \int_{0}^{1} \wNet_{t,i}  di
\end{aligned}\end{gathered}\end{equation}
while per capita assets are aggregate assets divided by aggregate population,
\begin{equation}\begin{gathered}\begin{aligned}
  \wNet_{t} & =  \WNet_{t}/1.
\end{aligned}\end{gathered}\end{equation}

Similarly, normalizing the population of firms to one yields
\begin{equation}\begin{gathered}\begin{aligned}
  \kap_{t} & =  K_{t}/1.
\end{aligned}\end{gathered}\end{equation}

Up to this point, we have allowed for the possibility that different
households might have different amounts of net worth.  We now impose
the assumption that every household is identical to every other
household.  This assumption rules out the presence of any debt in
equilibrium (if all households are identical, they cannot all be in
debt - who would they owe the money to?).  Indeed, in this case, the
aggregate capital stock per capita will equal the aggregate level of
net worth, $\kap_{t}=\wNet_{t}$.\footnote{The results do not change if we
  permit differences in the levels of wealth across households, but
  this is because we are assuming CRRA utility, perfect certainty,
  perfect capital markets, and various other things.  When any of
  these assumptions is relaxed, the distribution of assets does
  matter.  For exploration of this more complex and realistic
  framework, see \cite{carroll:brookings}, \cite{aiyagari:ge},
  \cite{ksHetero}, \cite{carrollRequiem}.  }

Thus, households' expectations about $\Wage_{t}$ and
$\rfree_{t}$ determine their saving decisions, which in turn determine the
aggregate path of $\kap_{t}$.

There is one important subtlety here, however.  In writing the 
consumer's budget constraint, we designated $\rfree_{t}$ as the net amount 
of income that would be generated by owning one more unit of net worth 
(e.g.\  capital).  But if we have depreciation of the capital stock, 
the net return to capital will be equal to the marginal product {\it minus}
depreciation.  The discussion of the firm's optimization problem did 
not consider depreciation because the firms do not own any capital; 
instead, they make a payment $\hat{\rfree}_{t}$ to the households for the 
privilege of using the households' capital.  Thus the net increment to 
a household's wealth if the household holds one more unit of capital 
will be
\begin{equation*}\begin{gathered}\begin{aligned}
        \rfree_{t} & =  \hat{\rfree}_{t}-\depr.
\end{aligned}\end{gathered}\end{equation*}

There is no depreciation of labor, so the labor market equilibrium
will be
\begin{equation*}\begin{gathered}\begin{aligned}
        \Wage_{t} & =  \hat{\Wage}_{t}
\\        & =  \fFunc(\kap_{t})- \hat{\rfree}_{t}\kap_{t}      
.
\end{aligned}\end{gathered}\end{equation*}

\section{The Perfect Foresight Equilibrium}

We assume that every household knows the aggregate production 
function, and understands the behavior of all the other households and 
firms in the economy.  Understanding all of this, suppose that 
households have some set of beliefs about the future path of the 
aggregate capital stock per capita $\{\kap_{t}\}_{t=0}^{\infty}$.  This 
belief about $\kap_{t}$ will imply beliefs about wages and interest rates 
as well $\{\Wage_{t},\rfree_{t}\}_{t=0}^{\infty}$.

The final assumption is that the equilibrium that comes about in this 
economy is the ``perfect foresight equilibrium.''  That is, consumers 
have the sets of beliefs such that, if they have those beliefs and act 
upon them, the actual outcome turns out to match the beliefs.  \opt{MarginNotes}{\marginpar{\tiny It can be shown that in this model there is only one `self-validating' set of beliefs.  Some other models may have multiple sets of beliefs all of which satisfy the condition that if everybody shares those beliefs, they will come true.  Interesting area of macro that Thomas has done some important work in.}}

Note now that using the fact that $\wNet_{t}=\kap_{t}$ in the perfect foresight
equilibrium we can rewrite the household's budget constraint as
\begin{equation}\begin{gathered}\begin{aligned}
        \dot{\kap}_{t} & =  \rfree_{t}\kap_{t} + \Wage_{t} - \cons_{t} %\label{eq:aggkdot}
\\  & =  (\hat{\rfree}_{t}-\depr)\kap_{t} + \Wage_{t} - \cons_{t} \label{eq:kdotind}
.
\end{aligned}\end{gathered}\end{equation}

Reproducing from \eqref{eq:cdotindiv},
\begin{equation}\begin{gathered}\begin{aligned}
        \dot{\cons}_{t}/\cons_{t} & =  \CRRA^{-1}(\rfree_{t}-\timeRate) %\label{eq:cdotagg}
\\  & =  \CRRA^{-1}(\hat{\rfree}_{t}-\depr-\timeRate)   \label{eq:cdotind}.
\end{aligned}\end{gathered}\end{equation}

Now compare these to the equations derived for the social planner's 
problem (with population growth and productivity growth zero) in a 
previous handout:
\begin{equation}\begin{gathered}\begin{aligned}
        \dot{\kap}_{t} & =  \fFunc(\kap_{t})-\depr \kap_{t}-\cons_{t}  \\
         & =  \hat{\rfree}_{t}\kap_{t}+\Wage_{t}-\depr \kap_{t}-\cons_{t}  \\
         & =  (\hat{\rfree}_{t}-\depr) \kap_{t}+\Wage_{t}-\cons_{t} \label{eq:kdotsp}
\end{aligned}\end{gathered}\end{equation}
and
\begin{equation}\begin{gathered}\begin{aligned}
        \dot{\cons}_{t}/\cons_{t} & =  \CRRA^{-1}(\fFunc^{\prime}(\kap_{t})-\depr-\timeRate). \label{eq:cdotsp}
\end{aligned}\end{gathered}\end{equation}

Since the equilibrium value of $\hat{\rfree}_{t}=\fFunc^{\prime}(\kap_{t})$, 
\eqref{eq:cdotsp} = \eqref{eq:cdotind}.  And \eqref{eq:kdotind} is 
identical to \eqref{eq:kdotsp}.  Thus, aggregate behavior of this economy 
is identical to the behavior of the social planner's economy!

This is a very convenient result, because it means that if we are 
careful about the exact assumptions we make we can often solve a 
social planner's problem and then assume that the solution also 
represents the results that would obtain in a decentralized economy.  
The social planner's solution and the decentralized solution are the 
same because they are maximizing the same utility function with 
respect to the same factor prices ($\rfree_{t}$ and $\Wage_{t}$).

When will the decentralized solution {\it not} match the social 
planner's solution?  One important case is when there are 
externalities in the behavior of individual households; another 
possible case is where there is idiosyncratic risk but no aggregate 
risk; basically, whenever the household's budget constraint or utility 
function differs in the right ways from the aggregate budget 
constraint or the social planner's preferences, there can be a 
divergence between the two solutions.

 
\input handoutBibMake

\end{document}

\begin{equation}\begin{gathered}\begin{aligned}
        \int_{0}^{\infty} \uFunc(c) e^{-\timeRate t}
\end{aligned}\end{gathered}\end{equation}
subject to
\begin{equation}\begin{gathered}\begin{aligned}
        \dot{\kap} & =  \fFunc(\kap) - \cons - \tau.
\end{aligned}\end{gathered}\end{equation}
which has Hamiltonian representation
\begin{equation}
        \Ham(\kap,c,\lambda) = \uFunc(c) + \lambda(\kap^{\alpha} - \cons - \tau)
\end{equation}

The first Hamiltonian optimization condition requires $\partial 
H/\partial \cons = 0$:
\begin{equation}\begin{gathered}\begin{aligned}
        \cons^{-\CRRA} & =  \lambda  \label{eq:H1A} \\
        -\CRRA \cons^{-\CRRA-1}\dot{\cons} & =  \dot{\lambda} %\label{eq:H1B}
\end{aligned}\end{gathered}\end{equation}

The second Hamiltonian optimization condition requires:
\begin{equation}\begin{gathered}\begin{aligned}
        \dot{\lambda} & =  \timeRate\lambda - \lambda(\partial \Ham/\partial k)  \\
         & =  \timeRate\lambda - \lambda \fFunc^{\prime}(\kap)
\\      \dot{\lambda}/{\lambda} & =  (\timeRate-\fFunc^{\prime}(\kap))
\\  \dot{\cons}/\cons  & =  \CRRA^{-1}(\fFunc^{\prime}(\kap)-\timeRate).
\end{aligned}\end{gathered}\end{equation}

Thus, the $\dot{\cons}=0$ locus in the phase diagram is unchanged.  
However, the $\dot{\kap}$ locus is shifted down by amount $\tau = 
\sigma$.

Now what happens if the government does not face a balanced budget
requirement?  Specifically, suppose $\debt$ is the level of government
debt, and the government's Dynamic Budget Constraint is \opt{MarginNotes}{\marginpar{\tiny LHS: Change in debt; RHS: spending minus taxes minus debt interest.}}
\begin{equation}\begin{gathered}\begin{aligned}
        \dot{d} & =  \sigma-\tau+r d
\end{aligned}\end{gathered}\end{equation}

The government's IBC will be the integral of the DBC:
\begin{equation}\begin{gathered}\begin{aligned}
        \debt_{0}+\int_{0}^{\infty}\sigma R   & =  \int_{0}^{\infty} \tau R   \\
        \debt_{0}+G_{0} & =  T_{0}
\end{aligned}\end{gathered}\end{equation}

The DBC of the representative family also changes.  They can now own 
either capital $\kap$ or government debt $\debt$.  If the family is to be 
indifferent between the two forms of assets, the interest rate must be 
the same.

\begin{equation}\begin{gathered}\begin{aligned}
        \cons + \dot{\wNet} & =  \Wage + \rfree \wNet - \tau  \\
        \wNet & =  \kap + \debt
\end{aligned}\end{gathered}\end{equation}

The assumption of labor augmenting technological progress was made 
because it implies that in steady-state $\dot{\cons}=0$ and $\dot{C}/C = 
\dot{Y}/Y = \dot{K}/K = \wGro$.

$\dot{\cons}/\cons = 0$ implies that at the steady-state value of $\bar{\kap}$
\begin{equation}\begin{gathered}\begin{aligned}
        \fFunc^{\prime}(\bar{\kap}) & =  \timeRate+\popGro+\depr+\CRRA \wGro  \\
        \alpha \bar{\kap}^{\alpha-1} & =  \timeRate+\popGro+\depr+\CRRA \wGro  \\
        \bar{\kap} & =  \left(\frac{\timePref+\popGro+\depr+\CRRA \wGro}{\alpha}\right)^{\frac{1}{\alpha-1}}
\\      \bar{\kap} & =  \left(\frac{\alpha}{\timePref+\popGro+\depr+\CRRA \wGro}\right)^{\frac{1}{1-\alpha}}   
\end{aligned}\end{gathered}\end{equation}

Thus, the steady-state capital/output ratio will be higher if capital 
is more productive ($\alpha$ is higher), and will be lower if 
consumers are more impatient, population growth is faster, 
depreciation is higher, or technological progress is higher (assuming 
$\CRRA>1$).

\end{document}
